\section{Conclusion}
\label{sec:conclusion}

The conclusion of saturation is unpleasant, but difficult to avoid. At the current resolution, the remaining errors in DHCD are driven more by label noise and handwriting ambiguities than by the model's ability. In other words, throwing a bigger recognizer at the problem does not buy a better recognizer. At that point, compactness is no longer a compromise. That will be the correct operating point.

\modelname{} reaches the accuracy of \textbf{99.73\,\%} with just the \textbf{1.11\,M} parameter and achieves statistical equivalence with the leading large-scale model baselines (McNemar $p = 0.345$) even before the advent of knowledge distillation. This means a parameter reduction of 15.6$\times$ and a CPU latency reduction of 9.5$\times$ without sacrificing accuracy. Why can't we move forward? The same \textbf{11-error} intrinsic floor breaks all tested configurations (including large-scale teacher ensembles), and ablation closes off the obvious escape route. Distillation target cleanliness, ensemble size, and test-time augmentation all have no statistically significant benefit. The floor is the ceiling.

The boundaries between the two must be clearly marked. The saturation claim is tied to certain conditions of DHCD, namely $32{\times}32$ resolution, concentrated label noise, and its particular pool of writers. Cleaner or higher resolution benchmarks may reveal capacity differences that DHCD may not reveal. The transfer results only cover \textbf{digit subset} (10 out of 46 classes). Consonant-level transferability is out of scope here, as there are no independent consonant target sets that share the label taxonomy of DHCD.

The practical implication for Devanagari recognition is to measure progress more honestly before claiming it. Without \emph{paired evaluation} -- McNemar or equivalent -- for a particular test case where the compared models also fail, a gain of a few hundredths of a percent is meaningless. The benchmark problem is even more difficult. DHCD served as a discriminator. The next necessary step is high-resolution scanning or a completely independent pool of writers.